\begin{longtable}{p{3.2cm}p{6.4cm}p{5.4cm}}
\toprule
\textbf{Tema} & \textbf{Decisão} & \textbf{Porquê} \\
\midrule
\endhead
Endereços & As páginas são servidas sem ``.html'' nem barra final (ex.: /advogados); o inglês fica em /en. & É o formato que a Cloudflare Pages serve sem redireções e mantém os endereços curtos. \\
Conteúdo & O conteúdo vive em ficheiros JSON validados (Zod) com campos PT/EN lado a lado, organizados como os futuros modelos do Sanity. & Permite testar tudo automaticamente e ligar o Sanity mais tarde sem mudar o desenho. \\
Conteúdo & Textos longos (artigos, páginas legais) guardados como blocos simples (parágrafo, título, lista). & Corresponde ao formato de texto rico do Sanity e evita dependências extra. \\
Conteúdo & Moradas completas dos escritórios marcadas ``a confirmar''; biografias limitadas a factos conhecidos. & Não inventar informação sobre o escritório nem sobre as pessoas. \\
Vídeo & Compressão com qualidade constante e teto de débito; o MP4 vem antes do WebM. & A textura da água fazia o WebM ultrapassar o limite de 2,5 MB; com o teto, o MP4 ficou mais leve. \\
Vídeo & Imagem de capa própria para telemóvel (recorte vertical) e capas de reserva limitadas a 100 KB. & A capa tem de coincidir com o vídeo vertical e aparecer depressa em 4G. \\
Vídeo & O botão de pausa do vídeo deixou de estar visível (pedido do cliente); aparece apenas quando alcançado com a tecla Tab. & Mantém a conformidade com WCAG 2.2.2 (forma de parar conteúdo em movimento) sem ocupar espaço visual. \\
Imagens & Imagens de banco ligeiramente dessaturadas por CSS. & Integram-se na paleta calma sem reprocessar ficheiros que vão ser substituídos. \\
Navegação & O menu completo aparece a partir de 1280 px; abaixo disso usa-se o botão Menu e a barra fixa Ligar/Contactar. & A 1024 px o menu completo com o botão não cabia sem partir linhas. \\
Rodapé & Rodapé claro (areia), como na referência Alden \& Vale. & Mantém o tom calmo; o verde escuro fica reservado ao topo e às chamadas para ação. \\
Situações & Separadores (ARIA tabs) em ecrãs largos e acordeão no telemóvel; sem JavaScript tudo fica visível. & Cada formato funciona melhor no seu ecrã e nada se perde sem JavaScript. \\
Formulário & Escritório preferido escolhido com botões (``chips''), por omissão ``Indiferente''. & Menos toques no telemóvel. \\
Formulário & Aviso ``sem JavaScript'' mostrado por CSS (@media scripting: none) em vez de ser removido por script. & Remover o aviso depois de carregar fazia a página saltar (CLS 0,106). \\
Mapas & Mapa ilustrado (SVG) desenhado no próprio site até ao clique; só então carrega o Google Maps. & Nenhum pedido à Google nem cookies antes da ação do visitante. \\
Coleções vazias & Casos, testemunhos e destaques sem conteúdo desaparecem da página inicial, do menu e do mapa do site. & Requisito do documento; nesta fase aparecem com o selo ``Exemplo ilustrativo''. \\
Segurança & Política de segurança de conteúdo (CSP) só por cabeçalho HTTP, sem 'unsafe-inline' nem 'unsafe-eval'. & O site não tem scripts nem estilos embutidos, por isso a política mais restrita funciona. \\
Segurança & A biblioteca de validação (Zod) corre em modo ``jitless''. & Os testes mostraram que ela tentava usar eval, que a CSP bloqueia — e bem. \\
Segurança & Auditoria de dependências bloqueia apenas dependências de produção (0 vulnerabilidades). & As 6 vulnerabilidades moderadas estão só em ferramentas de teste usadas localmente. \\
Desempenho & Pré-carregamento da letra DM Sans; a letra do logótipo carrega normalmente. & Evita saltos de texto (CLS) sem roubar largura de banda à imagem de capa em 4G. \\
Testes & Lighthouse com pedido de aquecimento e mediana de 3 execuções por página. & Execuções únicas a frio davam valores anómalos (2,9 s) que não representam o site. \\
Testes & Playwright e Lighthouse usam o Google Chrome já instalado. & Evita descarregar 150 MB de browsers de teste. \\
Testes & Na versão inglesa, nomes próprios (``Sociedade de Advogados'', ``Ordem dos Advogados'', linha do logótipo) não contam como texto português. & São nomes, não texto de interface por traduzir. \\
Telemóvel & Os botões Ligar/Contactar passaram de barra em largura total a dois botões redondos no canto inferior direito. & Pedido do cliente: menos chamativos e sem ocupar a parte de baixo do ecrã. \\
Páginas secundárias & Topo de imagem em largura total, do topo da página, com fio bronze no fundo; faixa editorial com fotografia ou bloco verde com etiqueta vertical e moldura deslocada. & Pedido do cliente (topos mais bonitos), sem copiar o estilo da referência e sem margens brancas no topo e nos lados. \\
Advogados & Cargo e nome escritos dentro da fotografia, sobre gradiente escuro. & Pedido do cliente, inspirado na referência. \\
20 anos & Nova secção na página inicial com o troféu do 20.º aniversário; o troféu é também a imagem do topo da História. & Pedido do cliente; o texto usa só factos conhecidos (fundação em 2002, 20 anos em 2022). \\
Telemóvel & O menu passou a gaveta que abre da esquerda, com uma margem de 3,5 rem à direita onde o site continua visível; tocar nessa margem fecha o menu. & Pedido do cliente. \\
Imagens & As caixas de imagem com proporção fixa cortam o excesso (overflow: hidden). & Sem isso, a fotografia (3:2) fazia a caixa (16:9) crescer quando carregava durante o scroll no telemóvel — problema reportado pelo cliente nos Destaques. \\
\bottomrule
\end{longtable}
